\section{Estado actual y trabajo pendiente}
\label{sec:estado}

\textbf{Avance.} El equipo estima completado alrededor del 80\,\% del
proyecto. Están terminados el núcleo algorítmico (las cuatro estructuras y
los siete requerimientos funcionales), la interfaz de consola, la suite de
pruebas y una primera medición de desempeño. El 20\,\% restante corresponde
a las funcionalidades y análisis de la entrega final que se listan abajo.

\textbf{Funcionalidades pendientes para la entrega final.}
\begin{itemize}
  \item \emph{Reclasificar la urgencia} de una solicitud pendiente
    (RF8). Con el localizador, retirarla de su nivel cuesta $O(1)$, pero
    insertarla en el nuevo nivel respetando su orden de llegada exige
    encontrar su posición, lo que en una lista cuesta $O(n_u)$.
  \item \emph{Carga por lotes} de solicitudes desde un archivo de texto,
    para reproducir escenarios de demostración y de medición.
  \item \emph{Estadísticas de atención}: tiempo de espera promedio por nivel
    y número de atenciones deshechas.
\end{itemize}

\textbf{Problemas técnicos abiertos.}
\begin{itemize}
  \item Cómo reclasificar en $O(1)$ sin perder el orden de llegada. Una
    alternativa en estudio es mantener en cada nivel una estructura
    ordenada por llegada solo para las solicitudes reclasificadas.
  \item El benchmark usa urgencias uniformes; falta medir distribuciones
    sesgadas (por ejemplo, mayoría de urgencias 4 y 5, como en la práctica)
    y medir en más de un equipo.
\end{itemize}

\textbf{Actividades para la entrega final.} Implementar y probar RF8 y la
carga por lotes; ampliar el benchmark con distribuciones de urgencia
sesgadas y con el peor caso de colisiones del índice; incorporar los
resultados y la retroalimentación de esta entrega al documento; y preparar
la demostración final.
